% 原书分章PDF第14—17页：第15.4节，完整聚集方差推导与模拟。
\section{方差估计与聚集性}

我们没有假定个体在统计上彼此独立，只假定各簇独立同分布。因此，观测单位是医院，上述 TMLE 的方差估计量需要考虑同一簇内个体之间的统计相依性。如果忽略聚集性，就会错误地认为簇内个体独立分布，从而低估真实的变异程度。

\subsection{聚集性与个体间干扰的区别}

聚集性的概念可能与个体间干扰（interference）混淆；因果分析通常假定不存在个体间干扰（Rubin，1980；Hudgens 和 Halloran，2008）。个体间干扰意味着，一名研究对象的潜在结局不依赖其他研究对象的暴露。在这里，这对应于以下假设：给定第二名婴儿的母乳喂养状态，第一名婴儿的母乳喂养不会影响第二名婴儿住院的概率。

相比之下，聚集性违反的是个体层面数据 $(O_{mi}:m,i)$ 独立同分布抽样的假设。簇内观测数据，包括暴露和结局，可能相关。不过，在不存在个体间干扰时，假定这种相关来自簇内人群的相似性，而各簇之间的人群则存在差异。在某些情形中，例如结局为传染病时，医院中心或社区等簇的内部可能存在个体间干扰。但我们认为 PROBIT 例子并非如此。

\subsection{使用 EIF 进行估计}

如果簇的数量足够多，就可以合理地使用（有效）影响函数估计方差，同时考虑有限个已知簇。以 $D^*(O_i)$ 表示研究对象 $i$ 的 EIF 值。将 $M$ 个簇记为 $Z_m$（$m=1,\ldots,M$），其中 $Z_m$ 表示属于第 $m$ 个簇的研究对象集合。参数 $\psi_0^{\bar a}$ 的 LTMLE 估计量记为 $\psi_n^{\bar a}$。尽管存在相依性，随着样本量增加，上述独立同分布型 LTMLE 应当仍然是一个以 $D^*(O)$ 为影响曲线的渐近线性估计量。因此，仍有
\[
\begin{aligned}
 \psi_n^{\bar a}-\psi_0^{\bar a}
 &\approx\frac{1}{n}\sum_{i=1}^{n}D^*(O_i)\\
 &=\frac{1}{n}\sum_{m=1}^{M}\sum_{i\in Z_m}D^*(O_i),
\end{aligned}
\]
其中，我们按照簇归属重新编排了单位的下标。再同时乘以和除以 $M$，得到
\[
 \psi_n^{\bar a}-\psi_0^{\bar a}
 =\frac{1}{M}\sum_{m=1}^{M}\sum_{i\in Z_m}D^*(O_i)\frac{M}{n}.
\]
由于簇与簇之间独立，可以将簇视为实验单位，其 EIF 为
$\sum_{i\in Z_m}D^*(O_i)\frac{M}{n}$。
所以，当 $M$ 足够大时，可以用各簇 EIF 的方差近似估计量的方差。

估计方差时，有
\[
\begin{aligned}
 \Var(\psi_n^{\bar a})
 &\approx\Var\left(\frac{1}{n}\sum_{i=1}^{n}D^*(O_i)\right)\\
 &=\frac{1}{n^2}\Var\left(\sum_{m=1}^{M}\sum_{i\in Z_m}D^*(O_i)\right).
\end{aligned}
\]
由影响函数的定义，$\EE(D^*(O_i))=0$，因此
$\Var(D^*(O_i))=\EE(D^*(O_i)^2)$。
此外，对同一簇内的两个单位 $i$ 和 $j$，有
$\Var(D^*(O_i)\times D^*(O_j))=\EE(D^*(O_i)\times D^*(O_j))$。
所以，上式等于
\[
 \frac{1}{n^2}\sum_{m=1}^{M}\sum_{i,j\in Z_m}
 \EE(D^*(O_i)\times D^*(O_j))I(i\ne j)
 +\EE(D(O_i^*)^2)I(i=j),
\]
其中，$I(\cdot)$ 是以逻辑陈述为自变量的指示函数。如果假定同一簇内各研究对象影响函数的协方差是簇特异的常数，即
$\rho_m=\EE(D^*(O_i)\times D^*(O_j))$，其中 $i\ne j$、$i,j\in Z_m$，
并且各研究对象的簇内方差 $\sigma_m^2=\EE(D^*(O_i)^2)$（$i\in Z_m$）在簇内也为常数，则上式简化为
\[
 \frac{1}{n^2}\sum_{m=1}^{M}n_m(n_m-1)\rho_m+n_m\sigma_m^2,
\]
其中，$n_m$ 是第 $m$ 个簇中的研究对象数。$\rho_m$ 和 $\sigma_m^2$ 的值，可以分别用各簇内部 EIF 的协方差和方差进行经验估计。

\subsection{模拟研究}

为考察考虑聚集性的重要性，我们开展了一项模拟研究。数据生成为 PROBIT 数据集的简化版本，与 Schnitzer 等（2014）中的模拟研究相似。在31个簇中，每个簇生成500名研究对象，总样本量为 $n=15{,}500$。基线协变量 $L_0=\{W,U\}$ 生成为彼此独立的 Gaussian 变量，其簇特异均值分别从不同的 Gaussian 分布中抽取。给定研究对象最近的历史和基线变量后，为每名研究对象独立生成时变变量 $(C_1,L_1,A_1,C_2,L_2,A_2,C_3,L_3)$。二元变量 $A_t$（$t=1,2$）表示继续母乳喂养；$C_t$（$t=1,2,3$）是删失指示变量；$L_t$（$t=1,2,3$）表示前一个时间区间内发生住院。结局 $Y=\sum_{t=1}^{3}L_t$ 是计数变量。基线变量 $U$ 仅是住院的风险因素，并不影响删失或暴露。

为了与真实 PROBIT 数据中观察到的关联相对应，我们特意设定：如果当前时点记录到住院，继续母乳喂养的可能性便降低；如果前一时点仍在母乳喂养，删失的可能性较低，而如果前一时点发生住院，删失的可能性较高。住院情况在给定基线变量和此前两次访视的母乳喂养状态后生成，使母乳喂养持续越久，住院概率越低。

我们分别在 $\bar a=(0,0)$ 和 $\bar a=(1,1)$ 下估计参数 $\psi^{\bar a}=\EE(Y^{\bar a})$。所关注的参数为
$\delta=\psi_0^{(1,1)}-\psi_0^{(0,0)}$，它对应于 PROBIT 研究中第一个目标参数。我们生成1000个数据集，使用修改后的序贯 LTMLE 估计目标参数，并以逻辑回归模型估计删失、暴露及住院概率。分析中将 $U$ 视为未测量变量。表15.4比较了以下三种标准误估计方法：（1）不考虑聚集性的 EIF 方差估计量；（2）考虑聚集性的 EIF 方差估计；（3）“成对”簇非参数自助法（Cameron 等，2008），它重抽的是簇，而非研究对象。

\begin{table}[htbp]
\centering
\caption{模拟研究：纳入与忽略聚集性时 LTMLE 方差估计的比较}
\label{tab:ch15-simulation}
\small
\begin{tabular}{lcc}
\toprule
方法 & 估计标准误 & 95\%置信区间覆盖率\\
\midrule
EIF，不考虑聚集性 & 0.007 & 85\\
EIF，考虑聚集性 & 0.009 & 92\\
自助法，考虑聚集性 & 0.009 & 98\\
\bottomrule
\multicolumn{3}{c}{真实效应 $=-0.030$，平均偏差 $=0\%$}\\
\end{tabular}
\end{table}

为估计标准误，我们对31个簇进行有放回重抽样，用重抽样数据计算估计值，重复200次，再由这200个估计值计算标准误。置信区间取重抽样估计值的第2.5和第97.5百分位数。表15.4表明，不考虑聚集性的方差估计会低估标准误，导致覆盖率不理想。考虑聚集性的 EIF 方法和自助法都得到接近理想的覆盖率，其中自助法略有过度覆盖，EIF 方法的覆盖略显不足。所有模拟和建模均使用 R 完成（R Development Core Team，2016）。Schnitzer 等（2014）进行了广泛的模拟研究，比较聚集数据情形下不同估计量的表现，包括使用参数模型的 LTMLE、使用超级学习器的 LTMLE、IPW、g 计算和序贯 g 计算。
\FloatBarrier
